\documentclass[aspectratio=169]{beamer}

\usetheme{Madrid}
\usecolortheme{default}
\setbeamertemplate{navigation symbols}{}

\usepackage[T1]{fontenc}
\usepackage[utf8]{inputenc}
\usepackage[italian]{babel}
\usepackage{listings}
\usepackage{xcolor}

\definecolor{codebg}{RGB}{245,247,250}
\definecolor{coderule}{RGB}{210,214,220}
\definecolor{codekw}{RGB}{0,72,130}
\definecolor{codecomment}{RGB}{90,90,90}

\lstdefinelanguage{KVProtocol}{
  morekeywords={PING,SET,GET,DELETE,EXISTS,KEYS,QUIT,OK,ERR,NOT_FOUND},
  sensitive=true,
  morecomment=[l]{\#}
}

\lstset{
  language=KVProtocol,
  basicstyle=\ttfamily\small,
  keywordstyle=\color{codekw}\bfseries,
  commentstyle=\color{codecomment},
  backgroundcolor=\color{codebg},
  frame=single,
  rulecolor=\color{coderule},
  showstringspaces=false,
  columns=fullflexible
}

\title{Key-Value Store v0}
\subtitle{Versione studenti}
\author{Architetture dei Sistemi Distribuiti}
\date{Lezione 09}

\begin{document}

\begin{frame}
  \titlepage
\end{frame}

\begin{frame}{Obiettivo della lezione}
  Alla fine di questa lezione dovreste saper distinguere chiaramente:

  \begin{itemize}
    \item interfaccia del servizio;
    \item contratto promesso al client;
    \item implementazione concreta del nodo;
    \item limiti che emergono appena il sistema cresce.
  \end{itemize}

  \vspace{0.5em}
  Punto chiave: la stessa API \texttt{GET/SET/DELETE} puo' avere semantiche molto diverse.
\end{frame}

\begin{frame}{Perche' partire da un KV store}
  Un key-value store sembra semplice, ma mette subito in evidenza problemi reali:

  \begin{itemize}
    \item modello dati minimale;
    \item operazioni intuitive;
    \item protocollo facile da testare;
    \item passaggio naturale da locale a distribuito;
    \item ottimo terreno per discutere correttezza e trade-off.
  \end{itemize}

  \vspace{0.5em}
  Non stiamo costruendo un database completo. Stiamo costruendo un caso di studio.
\end{frame}

\begin{frame}{Interfaccia vs implementazione}
  \begin{block}{Interfaccia}
    Che cosa puo' chiedere il client al servizio.
  \end{block}

  \begin{block}{Contratto}
    Quali garanzie riceve il client quando il servizio risponde.
  \end{block}

  \begin{block}{Implementazione}
    Come il server realizza quel contratto.
  \end{block}

  \vspace{0.5em}
  Errore classico: confondere la struttura dati locale con il servizio esposto in rete.
\end{frame}

\begin{frame}{Domanda guida}
  \begin{alertblock}{Domanda guida}
    Che cosa significa davvero che una \texttt{SET k v} e' andata a buon fine?
  \end{alertblock}

  Nella versione di oggi significa solo:

  \begin{itemize}
    \item il processo server ha aggiornato il proprio stato in RAM;
    \item il server ha inviato \texttt{OK} sulla connessione;
    \item nessuna persistenza e nessuna replica sono implicite.
  \end{itemize}

  Questa domanda guidera' tutte le prossime evoluzioni.
\end{frame}

\begin{frame}[fragile]{Protocollo v0}
  \begin{itemize}
    \item trasporto TCP;
    \item protocollo testuale UTF-8;
    \item una richiesta per riga;
    \item una risposta per riga.
  \end{itemize}

  \begin{lstlisting}
PING
OK PONG

SET corso asd
OK

GET corso
OK asd
  \end{lstlisting}
\end{frame}

\begin{frame}{Operazioni esposte}
  Il primo contratto del servizio include:

  \begin{itemize}
    \item \texttt{PING}: verifica di raggiungibilita';
    \item \texttt{SET key value}: scrittura;
    \item \texttt{GET key}: lettura;
    \item \texttt{DELETE key}: rimozione;
    \item \texttt{EXISTS key}: presenza della chiave;
    \item \texttt{KEYS}: elenco delle chiavi;
    \item \texttt{QUIT}: chiusura della sessione.
  \end{itemize}
\end{frame}

\begin{frame}{Scelte di contratto}
  Alcune decisioni sono piccole solo in apparenza:

  \begin{itemize}
    \item chiave assente: \texttt{NOT\_FOUND} oppure \texttt{ERR}?
    \item valore vuoto: va distinto da chiave assente?
    \item \texttt{DELETE} su chiave mancante: errore o no-op?
    \item \texttt{KEYS}: primitiva essenziale o comando diagnostico?
    \item ack di \texttt{SET}: successo locale o commit robusto?
  \end{itemize}

  Ogni scelta qui anticipa problemi che nel distribuito diventano costosi.
\end{frame}

\begin{frame}[fragile]{Casi limite da discutere}
  \begin{lstlisting}
GET corso
NOT_FOUND

SET corso
ERR usage: SET <key> <value>

SET corso ""
OK

GET corso
OK ""
  \end{lstlisting}

  Domanda: il protocollo distingue davvero tutti i casi che ci interessano?
\end{frame}

\begin{frame}{Che cosa promette davvero il v0}
  Garanzie minime:

  \begin{itemize}
    \item richieste della stessa connessione processate in ordine;
    \item mappa chiave-valore coerente nel singolo processo;
    \item \texttt{GET} dopo \texttt{SET} sullo stesso nodo vede il valore scritto.
  \end{itemize}

  Non garanzie:

  \begin{itemize}
    \item persistenza su disco;
    \item tolleranza ai guasti;
    \item replica;
    \item scalabilita';
    \item consistenza distribuita.
  \end{itemize}
\end{frame}

\begin{frame}{Laboratorio}
  Lavoro richiesto:

  \begin{enumerate}
    \item definire il contratto del protocollo;
    \item implementare il server TCP single-node;
    \item gestire errori e casi limite;
    \item testare manualmente il comportamento osservabile;
    \item giustificare almeno una scelta semantica.
  \end{enumerate}

\end{frame}

\begin{frame}{Domande da tenere aperte}
  Durante il laboratorio chiedetevi sempre:

  \begin{itemize}
    \item che cosa promette davvero \texttt{SET};
    \item come distinguere chiave assente e valore vuoto;
    \item quali errori sono di protocollo e quali applicativi;
    \item quali scelte locali oggi diventano problemi domani.
  \end{itemize}
\end{frame}

\begin{frame}{Perche' \texttt{KEYS} e' didatticamente interessante}
  Oggi \texttt{KEYS} sembra innocuo:

  \begin{itemize}
    \item stato locale;
    \item costo basso;
    \item semantica ovvia.
  \end{itemize}

  Domani diventa ambiguo:

  \begin{itemize}
    \item su quale nodo?
    \item su quali repliche?
    \item su quali shard?
    \item con quale consistenza?
  \end{itemize}
\end{frame}

\begin{frame}{Come evolve il percorso}
  Dopo il v0, la stessa interfaccia verra' messa sotto pressione:

  \begin{itemize}
    \item concorrenza e race condition;
    \item persistenza e crash recovery;
    \item replica primaria-secondaria;
    \item failover e leader election;
    \item quorum e conflitti;
    \item sharding e riequilibrio.
  \end{itemize}

  Ogni evoluzione dovra' aggiornare il contratto, non solo il codice.
\end{frame}

\begin{frame}{Messaggio finale}
  \begin{block}{Idea centrale}
    Nei sistemi distribuiti l'interfaccia non basta. Conta il significato osservabile delle risposte.
  \end{block}

  Se non sapete spiegare cosa promette oggi \texttt{OK}, domani non potrete spiegare:

  \begin{itemize}
    \item cosa e' consistente;
    \item cosa sopravvive a un crash;
    \item cosa succede durante un failover;
    \item quali trade-off state accettando.
  \end{itemize}
\end{frame}

\end{document}
